\chapter{Manual de Usuario del Recomendador CNI} \label{anexo:manual_usuario}

Este anexo describe cómo utilizar el prototipo recomendador ---la herramienta central del proyecto--- para obtener una recomendación de plugin CNI fundamentada en los datos del \emph{testbed}. El recomendador es una aplicación web de página única (SPA) que no requiere instalación por parte del usuario final: se accede desde un navegador a la versión publicada en GitHub Pages del repositorio de resultados (\texttt{K8S-CNI-Results}).

\section{Acceso a la herramienta}

La aplicación está disponible como sitio estático publicado desde la carpeta \texttt{docs/recommender} del repositorio \texttt{K8S-CNI-Results} mediante GitHub Pages. Basta con abrir la URL del sitio en cualquier navegador moderno; no se necesita usuario, contraseña ni componentes adicionales, ya que la SPA consume directamente el archivo de datos procesados \texttt{cni-data.json}.

\section{Flujo de uso paso a paso}

\begin{enumerate}
\item \textbf{Responder el cuestionario guiado.} La pantalla principal presenta cinco preguntas en lenguaje no técnico ---velocidad de respuesta, estabilidad de la conexión, volumen de datos, límite de recursos del servidor y nivel de seguridad requerido---, cada una con un deslizador (\emph{slider}) de 1 a 5. Un campo de texto opcional permite indicar para qué es el proyecto.
\item \textbf{Revisar el perfil calculado.} A partir de las respuestas, la herramienta deriva los factores de ponderación de las métricas (seguridad, latencia, throughput, estabilidad y recursos) y los normaliza automáticamente al 100\,\%. El panel lateral muestra el perfil resultante y el peso asignado a cada criterio. Los perfiles \emph{Fintech / Seguridad crítica}, \emph{Streaming / Alto rendimiento} e \emph{IoT / Recursos limitados} documentados en esta tesis corresponden a combinaciones de referencia de esos mismos deslizadores; la aplicación no los ofrece como opciones preseleccionables.
\item \textbf{Consultar la recomendación.} La aplicación calcula el ranking multicriterio (MCDA) para el perfil derivado y resalta el CNI con mayor puntaje como recomendación principal.
\item \textbf{Revisar la justificación.} La interfaz muestra las métricas brutas promedio de cada CNI ---throughput, latencia, retransmisiones, consumo de CPU y RAM, e impacto de las Network Policies--- en gráficas comparativas, de modo que la recomendación sea explicable y no una ``caja negra''.
\item \textbf{Obtener la orientación de despliegue.} Para el CNI recomendado, la herramienta presenta el comando de instalación del plugin y una política de denegación por defecto, ambos copiables al portapapeles. Estos elementos sirven de punto de partida para el despliegue documentado en el \cref{anexo:manual_instalacion}; el ajuste fino de cada plugin queda a cargo del operador.
\end{enumerate}

\section{Interpretación de los resultados}

El puntaje de cada CNI se expresa en una escala de utilidad de 1 a 5, donde 5 representa la mejor alternativa en la métrica correspondiente (\cref{eq:mcda-suma,eq:mcda-norm}). Conviene tener presente la \textbf{restricción no compensatoria de seguridad}: cuando el perfil exige un nivel de seguridad alto, un CNI que no aplique Network Policies en su plano de datos recibe una penalización que lo descarta como opción viable, con independencia de su rendimiento. Por ello, en perfiles como Fintech, un CNI muy eficiente pero sin soporte de políticas no aparecerá como recomendado.

\section{Resolución de problemas frecuentes}

\begin{itemize}
\item \emph{El indicador de la cabecera muestra ``Datos de referencia''.} La aplicación no pudo descargar \texttt{cni-data.json} y está usando el conjunto de métricas embebido. Verifique la conexión a Internet y que el archivo esté publicado ---se regenera con \texttt{npm run sync:data} (\cref{anexo:manual_instalacion})--- y recargue la página, ya que la descarga no se reintenta automáticamente.
\item \emph{Los puntajes apenas cambian al ajustar los deslizadores.} Los deslizadores solo alteran el peso relativo de cada criterio: si todos se mueven en la misma dirección, el vector normalizado varía poco. Acentúe la diferencia entre las prioridades que sí importan para el despliegue.
\item \emph{Quiero evaluar un CNI nuevo o datos más recientes.} La herramienta está diseñada para mejora continua: al incorporar nuevas corridas en el repositorio de resultados y regenerar los datos, el ranking se recalcula sin modificar la aplicación.
\end{itemize}
